\subsection{A smooth-total-space obstruction for the proposed finite flat torsors}
\paragraph{A2Gs (Construction exclusion; submitted for independent audit).}
Let $Z$ be a smooth projective surface over an algebraically closed
field, let $D\subset Z$ be a smooth ample integral curve, and suppose
there is a line subbundle
\[
 \mathcal F\simeq\mathcal O_Z(-D)\ \subset\ T_Z
\]
whose restriction to $D$ is tangent to $D$. Here ``subbundle'' means
that the quotient is locally free. Such a triple does not exist.
No integrability assumption is needed for this assertion.

\paragraph{Chern and adjunction consequences.}
The inclusion on $D$ factors through $T_D$. It is nonzero on every
fibre, hence is an isomorphism of line bundles:
\[
 \mathcal O_D(-D)\simeq T_D.
\]
Taking degrees and using adjunction gives
\[
 D^2=2g(D)-2>0,\qquad K_Z\cdot D=0.
\]
The quotient line bundle in $0\to\mathcal F\to T_Z\to\mathcal Q\to0$
has first Chern class $-K_Z+D$. Consequently
\[
 c_2(Z)=(-D)\cdot(-K_Z+D)=-D^2<0.                 \tag{A2Gs.1}
\]
If any positive multiple of $K_Z$ has a nonzero section, its effective
divisor has zero intersection with the ample curve $D$, and is therefore
zero. Thus some positive multiple of $K_Z$ is trivial and
$\kappa(Z)=0$. Such $Z$ is minimal: a $(-1)$-curve would have
$K_Z\cdot E=-1$, contradicting torsion of $K_Z$. The classification
of minimal surfaces of Kodaira dimension zero gives
$c_2(Z)\in\{0,12,24\}$, contrary to (A2Gs.1). These classification
facts, including the characteristic-two and -three cases, are recorded
in \cite[Section 7, pp.~26--30]{DGJ_LiedtkeSurfaces2013}.
Hence $\kappa(Z)=-\infty$.

\paragraph{The ruled case, including infinitely near blowups.}
By \cite[Theorems 4.4--4.5, p.~21]{DGJ_LiedtkeSurfaces2013}, $Z$
is obtained by $r\geq0$ point blowups from a minimal ruled surface
$Z_0\to C$ (or a rational minimal surface). Since
$c_2(Z)=4(1-g(C))+r<0$ in the ruled case and rational surfaces have
positive $c_2$, the rational case is excluded and $g:=g(C)\geq2$.
Write $F$ for a fibre class, choose a section class $C_0$ with
$C_0^2=-e$, and use the total-transform exceptional classes $E_i$.
Numerically,
\[
 D=aC_0+bF-\sum_{i=1}^r m_iE_i,\qquad
 K_Z=-2C_0+(2g-2-e)F+\sum_{i=1}^r E_i.
\]
All intersection products here use pullbacks of $C_0,F$ to $Z$.
The integers satisfy
\[
 a=D\cdot F\geq1,\qquad 0<m_i<a.                 \tag{A2Gs.2}
\]
Indeed $E_i$ is a nonzero effective total-transform divisor, so
$m_i=D\cdot E_i>0$. If its blowup centre lies over $c\in C$, then
$F_c-E_i$ is also effective and nonzero: at the step creating $E_i$
the fibre has positive multiplicity at the centre, and subtracting
one exceptional divisor from its total transform leaves an effective
nonzero divisor. Subsequent pullbacks preserve this fact. Ampleness
then gives $a-m_i=D\cdot(F_c-E_i)>0$. This proves (A2Gs.2) also for
infinitely near centres.

The equality $K_Z\cdot D=0$ reads
\[
 2b=a(e+2g-2)+\sum_i m_i,
\]
and therefore
\[
 D^2=2a^2(g-1)+\sum_i m_i(a-m_i).
\]
Combining this with $c_2(Z)=4(1-g)+r=-D^2$ gives
\[
 (2a^2-4)(g-1)+\sum_i\bigl(m_i(a-m_i)+1\bigr)=0. \tag{A2Gs.3}
\]
If $a\geq2$, every term on the left is nonnegative and the first
is strictly positive. If $a=1$, (A2Gs.2) forces $r=0$, and (A2Gs.3)
becomes $-2(g-1)=0$. Both cases are impossible. This proves the assertion.

\paragraph{Application to the proposed group scheme.}
Let $X$ be a smooth projective surface in characteristic $p>0$,
$Y\subset X$ a smooth ample integral curve, $L=\mathcal O_X(Y)$,
and $s$ its canonical section. The rank-$p$ group scheme $G_s$ is
locally
\[
 \mathcal O_X[u]/(u^p),\qquad
 \Delta(u)=u\otimes1+1\otimes u+s\,u\otimes u.
\]
Our line convention is explicit: if local components of $s$ satisfy
$s_j=g_{ij}^{-1}s_i$, then $u_j=g_{ij}u_i$. Thus the algebra
generators transform as $L$, and $\operatorname{Lie}(G_s)=L^{-1}$.
The map $G_s\to\mu_p$ is $u\mapsto1+su$; it is an isomorphism
off $Y$, while the fibre on $Y$ is $\alpha_{L^{-1}|Y}$.

There is no $G_s$-torsor $\pi:Z\to X$ for which $Z$ is smooth and
$D=\pi^*Y$ is smooth integral. Indeed, $Z$ is projective because
$\pi$ is finite, and $D$ is ample. The action differentiates to
\[
 \pi^*\operatorname{Lie}(G_s)=\mathcal O_Z(-D)\longrightarrow T_Z.
\]
For each geometric point $z$, the orbit map from the fibre group
scheme to $Z$ is a closed immersion: the torsor identity identifies
it with the fibre of $\pi$ through $z$. Its tangent map is injective.
The displayed map is consequently a line subbundle, and its image
is tangent to $D$ because it annihilates pullbacks of functions from
$X$. The preceding assertion applies.

\paragraph{Exact limitation.}
This excludes the proposed construction only with a smooth total
space and a smooth full inverse image. A torsor with a singular
normalization is not covered: resolving it can make the pullback of
$Y$ merely nef and big and can introduce zeros or discrepancies in
the tangent distribution. Neither an arbitrary Picard point kernel
nor an arbitrary normalization of its $\mu_p$-torsor has been shown
to satisfy these extra hypotheses. The general A2 question remains open.
